%% terms-and-conditions.tex
%% Template per Termini e Condizioni — PoliNetwork APS
%%
%% Compila con: xelatex terms-and-conditions.tex
%% (dalla root del progetto)
%%
\documentclass[legal]{../core/polinetwork}

%% ====================================================================
%% DATI DOCUMENTO
%% ====================================================================

\pnTitle{Termini e Condizioni di Utilizzo}
\pnSubtitle{del servizio [Nome Servizio]}
\pnDate{1 aprile 2026}
\pnVersion{1.0}
\pnSubject{Termini e Condizioni --- [Nome Servizio]}

%% ====================================================================

\begin{document}
\thispagestyle{firstpage}

\pnMakeHeader
\pnMakeMetadata

%% ====================================================================
%% CONTENUTO
%% ====================================================================

I presenti Termini e Condizioni regolano l'utilizzo di [servizio/piattaforma]
fornito da \textbf{PoliNetwork APS}. L'accesso e l'utilizzo del servizio
implicano l'accettazione dei presenti termini.

%% --- Art. 1 ---
\pnArticle{Definizioni}

Ai fini dei presenti Termini e Condizioni si intende per:

\pnClause{%
  \textbf{``Associazione''}: PoliNetwork APS, associazione di promozione sociale
  senza scopo di lucro.
}

\pnClause{%
  \textbf{``Servizio''}: [descrizione del servizio offerto].
}

\pnClause{%
  \textbf{``Utente''}: qualsiasi persona fisica che accede al Servizio.
}

%% --- Art. 2 ---
\pnArticle{Oggetto}

\pnClause{%
  I presenti Termini e Condizioni disciplinano le modalità di accesso e
  di utilizzo del Servizio da parte degli Utenti.
}

\pnClause{%
  L'Associazione si riserva il diritto di modificare i presenti Termini 
  in qualsiasi momento, dandone comunicazione agli Utenti mediante 
  [specificare canale di comunicazione].
}

%% --- Art. 3 ---
\pnArticle{Accesso al Servizio}

\pnClause{%
  L'accesso al Servizio è gratuito e aperto a [specificare destinatari].
}

\pnClause{%
  L'Utente si impegna a fornire informazioni veritiere e aggiornate
  al momento della registrazione, ove richiesta.
}

%% --- Art. 4 ---
\pnArticle{Obblighi dell'Utente}

L'Utente si impegna a:

\begin{itemize}
  \item utilizzare il Servizio nel rispetto delle leggi vigenti;
  \item non compiere attività che possano compromettere la sicurezza
        del Servizio o dei dati degli altri Utenti;
  \item non diffondere contenuti illeciti, offensivi o contrari
        all'ordine pubblico o al buon costume.
\end{itemize}

%% --- Art. 5 ---
\pnArticle{Proprietà intellettuale}

\pnClause{%
  Tutti i contenuti presenti sul Servizio (testi, grafica, loghi, 
  icone, software) sono di proprietà dell'Associazione o dei rispettivi
  titolari dei diritti e sono protetti dalle leggi applicabili.
}

\pnClause{%
  È vietata qualsiasi forma di riproduzione, distribuzione o utilizzo
  non autorizzato dei contenuti del Servizio.
}

%% --- Art. 6 ---
\pnArticle{Limitazione di responsabilità}

\pnClause{%
  L'Associazione si impegna a garantire la continuità e la sicurezza
  del Servizio, ma non può essere ritenuta responsabile per
  interruzioni, malfunzionamenti o perdita di dati dovuti a cause
  di forza maggiore o a eventi al di fuori del proprio controllo.
}

\pnClause{%
  Il Servizio è fornito ``così com'è'' (\textit{as is}), senza garanzie
  di alcun tipo, esplicite o implicite.
}

%% --- Art. 7 ---
\pnArticle{Trattamento dei dati personali}

\pnClause{%
  Il trattamento dei dati personali degli Utenti è disciplinato
  dalla Informativa Privacy disponibile su
  \href{https://polinetwork.org}{polinetwork.org}.
}

%% --- Art. 8 ---
\pnArticle{Legge applicabile e foro competente}

\pnClause{%
  I presenti Termini e Condizioni sono regolati dalla legge italiana.
}

\pnClause{%
  Per qualsiasi controversia derivante dall'interpretazione o
  dall'esecuzione dei presenti Termini, le parti si impegnano a
  tentare una risoluzione in via amichevole. In caso di mancata
  composizione, sarà competente il Foro di Milano.
}

%% ====================================================================

\vfill
\pnContactBlock

\end{document}
